\documentclass[a4paper,10pt]{article}
\usepackage[utf8]{inputenc}
\usepackage[T1]{fontenc}
\usepackage[french]{babel}
\usepackage{amsmath,amssymb}
\usepackage{geometry}
\usepackage{array,tabularx}
\usepackage{tikz}
\geometry{margin=1.6cm}
\pagestyle{empty}
\newcommand{\ligne}{\par\vspace{0.55cm}\noindent\dotfill\par}
\newcommand{\carreaux}[1]{\begin{center}\begin{tikzpicture}\draw[step=0.5cm,gray!35,very thin] (0,0) grid (17,#1);\end{tikzpicture}\end{center}}
\newenvironment{exercice}[2]{\paragraph{Exercice #1 \hfill #2 points}\mbox{}\par}{\medskip}
\begin{document}
\begin{center}\Large\bfseries Devoir surveillé - Fonctions linéaires\end{center}
\vspace{2mm}
\begin{exercice}{1}{4}Donner la définition d'une fonction linéaire et la forme de son expression.\end{exercice}\begin{exercice}{2}{5}Soit $f(x)=-3x$. Calculer $f(4)$, $f(-5)$ et l'antécédent de $18$.\end{exercice}\begin{exercice}{3}{5}Une fonction linéaire vérifie $g(6)=15$. Déterminer son coefficient et calculer $g(10)$.\end{exercice}\begin{exercice}{4}{6}Représenter $h(x)=1,5x$ dans un repère et expliquer pourquoi sa courbe passe par l'origine.\carreaux{6}\end{exercice}\end{document}
